% =====================================================================
%  Nota suplementar à apostila de Relatividade Geral (UDESC/CCT)
%  Por que A = D na dedução da transformação de Lorentz (Seção 1.4)
%
%  Compile de dentro desta pasta com:
%     TEXINPUTS=../fonte: latexmk -xelatex nota_A_igual_D.tex
% =====================================================================
\documentclass[12pt,a4paper,oneside]{book}

\usepackage[brazil]{babel}
\usepackage{estilo-rg}

\usepackage{amsfonts}
\usepackage{array}
\usepackage{geometry}

\geometry{top=2.8cm,bottom=2.8cm,left=3.0cm,right=3.0cm,
          headheight=23pt,headsep=16pt}

\hypersetup{
  colorlinks=true,
  linkcolor=rgPetroleo,
  urlcolor=rgPetroClaro,
  pdftitle={Por que A = D: nota sobre a dedução da transformação de Lorentz},
  pdfauthor={UDESC/CCT}
}

% numeração de seção sem capítulo
\renewcommand{\thesection}{\arabic{section}}
\renewcommand{\theequation}{\arabic{equation}}

\markboth{Nota suplementar}{Transformação de Lorentz}

\begin{document}

% ---------------------------------------------------------------------
%  Cabeçalho da nota
% ---------------------------------------------------------------------
\thispagestyle{plain}

\noindent
{\sffamily\bfseries\footnotesize\color{rgAmbar}NOTA SUPLEMENTAR À SEÇÃO 1.4}

\vspace{4pt}
\noindent
{\sffamily\bfseries\LARGE\color{rgPetroleo}Por que $A=D$?}

\vspace{2pt}
\noindent
{\sffamily\large\color{rgGrafite}O papel da simetria $S\leftrightarrow S'$
na dedução da transformação de Lorentz}

\vspace{6pt}
\noindent\textcolor{rgLinha}{\rule{\linewidth}{0.8pt}}

\vspace{10pt}

\noindent
A caixa \textsc{Resultado Central} da p.~14 deduz a transformação de Lorentz
a partir da invariância do intervalo e, em certo ponto, afirma:

\begin{quote}\itshape\small\color{rgGrafite}
``A simetria entre os dois referenciais (o que $S$ vê como velocidade $v$,
$S'$ vê como $-v$, e nenhum referencial é privilegiado) impõe $A=D$.''
\end{quote}

\noindent
Esta nota responde a duas perguntas que a frase deixa em aberto:
\textbf{(i)} o que exatamente essa simetria tem a ver com a equação
$D^2v^2-A^2=-1$, que aparece imediatamente antes; e \textbf{(ii)} qual é,
de fato, o conteúdo lógico da hipótese --- que se revela bem menor do que
parece.

% =====================================================================
\section{O sistema completo, e o que a caixa usa dele}
% =====================================================================

A transformação postulada é linear,
\begin{equation}
  t' = At + Bx, \qquad x' = Ct + Dx,
  \label{eq:linear}
\end{equation}
com quatro incógnitas $A,B,C,D$, funções apenas de $v$. Sobre elas a caixa
impõe quatro equações:

\vspace{4pt}
\begin{center}
\small
\begin{tabular}{@{}l l l@{}}
\toprule
\textbf{Origem} & \textbf{Equação} & \textbf{Usada na caixa?}\\
\midrule
Condição 1 (cinemática) & $C=-Dv$ & sim \\
Condição 2, coef.\ de $\dd t^2$ \hspace{1em} & (E1)\quad $-A^2+C^2=-1$ & sim \\
Condição 2, coef.\ de $\dd x^2$ & (E2)\quad $-B^2+D^2=1$  & \textcolor{rgTerracota}{\textbf{não}} \\
Condição 2, coef.\ de $\dd t\,\dd x$ & (E3)\quad $-AB+CD=0$ & sim (no final) \\
\bottomrule
\end{tabular}
\end{center}
\vspace{4pt}

\noindent
O texto usa a Condição 1 e (E1) para escrever $D^2v^2-A^2=-1$; invoca então
a simetria para obter $A=D$; e só no fim recorre a (E3) para achar
$B=-\gamma v$. A equação (E2) \emph{nunca é utilizada}. É aí que mora toda a
questão.

\begin{atencao}
Antes de qualquer conta, uma observação de contagem. São
\textbf{quatro incógnitas} e \textbf{quatro equações}. Um sistema assim não
deixa nenhuma liberdade \emph{contínua} sobrando: o conjunto solução é
discreto. Logo, qualquer hipótese física adicional --- inclusive a simetria
--- não pode estar fixando um parâmetro livre. No máximo pode estar
\textbf{escolhendo um sinal}. E é exatamente isso que ela faz.
\end{atencao}

% =====================================================================
\section{Primeiro: o que a simetria realmente afirma}
% =====================================================================

Um esclarecimento antes da álgebra. A simetria \textbf{não age sobre} a
equação $D^2v^2-A^2=-1$. Ela é uma hipótese independente, cujo resultado é
\emph{substituído} naquela equação depois. As duas coisas aparecem na mesma
frase da caixa, o que sugere --- erradamente --- que uma esteja sendo
deduzida da outra.

Em fórmula, a simetria é a imagem espelhada da Condição 1. Aquela fixou como
a origem de $S'$ se move vista de $S$; esta fixa como a origem de $S$ se move
vista de $S'$. Basta pôr $x=0$ (a linha de mundo da origem de $S$) em
\eqref{eq:linear}:
\begin{equation}
  t' = At, \qquad x' = Ct = -Dvt,
\end{equation}
onde já se usou $C=-Dv$. A velocidade da origem de $S$ \emph{medida em $S'$}
é portanto
\begin{equation}
  \frac{x'}{t'} \;=\; -\frac{D}{A}\,v .
  \label{eq:reciproca}
\end{equation}
O postulado de que nenhum referencial é privilegiado exige que esse valor
seja exatamente $-v$: se $S$ vê $S'$ afastar-se com velocidade $v$, então
$S'$ deve ver $S$ afastar-se com a mesma velocidade, em sentido oposto.
Impondo isso em \eqref{eq:reciproca},
\begin{equation}
  -\frac{D}{A}\,v = -v \qquad\Longrightarrow\qquad A=D .
\end{equation}

\noindent
Esse é todo o argumento da caixa, agora explícito: uma segunda condição
\emph{cinemática}, do mesmo tipo da primeira. Sem ela, ficaria a esquisitice
de $S$ medir $v$ para $S'$ enquanto $S'$ mede, digamos, $-1{,}3\,v$ para $S$.
Só depois disso se volta a $D^2v^2-A^2=-1$ com $A=D\equiv\gamma$, obtendo
$\gamma^2(1-v^2)=1$.

% =====================================================================
\section{A conta que a caixa não faz}
% =====================================================================

Agora o ponto interessante: \textbf{a álgebra já entrega $D^2=A^2$ sozinha}.
Substitua $C=-Dv$ nas três equações da invariância:
\begin{align}
  \text{(E1$'$)}&\qquad A^2 - D^2v^2 = 1, \label{eq:E1}\\
  \text{(E2$'$)}&\qquad D^2 - B^2 = 1,    \label{eq:E2}\\
  \text{(E3$'$)}&\qquad -AB + (-Dv)D = 0 \;\Longrightarrow\; AB = -D^2v.
  \label{eq:E3}
\end{align}

\paragraph{Passo 0 --- pode-se dividir por $A$.}
De \eqref{eq:E1}, $A^2 = 1 + D^2v^2 \ge 1$, logo $|A|\ge 1$ e em particular
$A\neq 0$. Não é preciosismo: é o que legitima o passo seguinte.

\paragraph{Passo 1 --- isole $B$ em \eqref{eq:E3}.}
\begin{equation}
  B = -\frac{D^2 v}{A}.
\end{equation}

\paragraph{Passo 2 --- leve isso em \eqref{eq:E2},} a equação que a caixa
deixou de lado:
\begin{equation}
  D^2 - \frac{D^4v^2}{A^2} = 1 .
\end{equation}

\paragraph{Passo 3 --- multiplique por $A^2$ e fatore.}
\begin{equation}
  A^2D^2 - D^4v^2 = A^2
  \qquad\Longrightarrow\qquad
  D^2\bigl(A^2 - D^2v^2\bigr) = A^2 .
\end{equation}

\paragraph{Passo 4 --- reconheça o parêntese.}
O termo entre parênteses é precisamente o lado esquerdo de \eqref{eq:E1},
que vale $1$. Restam apenas

\begin{central}
\begin{equation}
  \boxed{\;D^2 = A^2\;}
\end{equation}
obtido \emph{sem} nenhuma menção a simetria, a referenciais privilegiados ou
à troca $S\leftrightarrow S'$. É puro encaixe entre as três equações da
invariância do intervalo mais a Condição 1.
\end{central}

% =====================================================================
\section{Os dois ramos, explicitamente}
% =====================================================================

De $D^2=A^2$ segue $D=\pm A$. Vale escrever os dois casos, porque o ramo
descartado não é nenhum monstro algébrico --- é uma transformação
perfeitamente legítima.

\paragraph{Ramo $D=+A$.} Então $C=-Av$ e $B=-A^2v/A=-Av$, de modo que
\begin{equation}
  t' = A(t-vx), \qquad x' = A(x-vt),
  \qquad \det = A^2 - A^2v^2 = +1 .
\end{equation}
Com \eqref{eq:E1}, $A^2(1-v^2)=1$; escolhendo $A>0$ obtém-se $A=\gamma$. É o
boost usual.

\paragraph{Ramo $D=-A$.} Então $C=+Av$ e $B=-A^2v/A=-Av$, de modo que
\begin{equation}
  t' = A(t-vx), \qquad x' = -A(x-vt),
  \qquad \det = -A^2 + A^2v^2 = -1 .
\end{equation}
É o \emph{mesmo boost composto com a reflexão} $x\to-x$. Ele satisfaz todas
as quatro equações, \textbf{inclusive a Condição 1}: a origem de $S'$
($x'=0$) continua sobre $x=vt$. O que ele viola é uma convenção que nunca
chegou a ser escrita --- a de que os eixos $x$ e $x'$ apontem para o mesmo
lado. Em $v=0$ ele dá $t'=t$, $x'=-x$: não tende à identidade.

% =====================================================================
\section{Leitura estrutural: isso tinha que acontecer}
% =====================================================================

O resultado $D^2=A^2$ não é uma coincidência algébrica. Calcule o
determinante da matriz da transformação,
$M=\begin{pmatrix} A & B\\ C & D\end{pmatrix}$, usando $C=-Dv$ e
$B=-D^2v/A$:
\begin{align}
  \det M &= AD - BC
         = AD - \Bigl(-\frac{D^2v}{A}\Bigr)(-Dv)
         = AD - \frac{D^3v^2}{A} \notag\\[4pt]
         &= \frac{D}{A}\underbrace{\bigl(A^2-D^2v^2\bigr)}_{=\,1}
         = \frac{D}{A}.
\end{align}
Ou seja, $\det M = D/A$, e portanto
\begin{equation}
  D^2=A^2 \iff \det M = \pm1,
  \qquad\qquad
  A = D \iff \det M = +1 .
\end{equation}

Mas $\det M=\pm1$ é automático para \emph{qualquer} transformação que
preserve uma forma quadrática. De $M^{\mathsf T}\eta M=\eta$, tomando
determinantes,
\begin{equation}
  (\det M)^2\,\det\eta = \det\eta
  \qquad\Longrightarrow\qquad
  (\det M)^2 = 1 .
\end{equation}
As equações (E1)--(E3) \emph{são} a afirmação $M^{\mathsf T}\eta M=\eta$: elas
caracterizam o grupo $O(1,1)$ inteiro, com suas quatro componentes
desconexas (identidade, $P$, $T$, $PT$). Não há como elas distinguirem $A=D$
de $A=-D$ --- isso é informação de \emph{componente}, não de equação. Nenhuma
manipulação algébrica poderia fornecê-la; é preciso um argumento físico, e é
esse o serviço prestado pela simetria.

% =====================================================================
\section{Conclusão}
% =====================================================================

\begin{central}
\textbf{(E2) e a hipótese de simetria são redundantes entre si.} A caixa faz
uma escolha: invoca a simetria e deixa (E2) ociosa. Poderia ter feito a
escolha oposta --- usar (E2) para obter $D^2=A^2$ e reservar a simetria
apenas para o sinal. O que não se pode é achar que a simetria esteja sendo
aplicada \emph{à} equação $D^2v^2-A^2=-1$: ela é um passo à parte, cujo
resultado é substituído ali depois.
\end{central}

Que (E2) fique satisfeita no fim é fácil de conferir: com $A=D=\gamma$ e
$B=-\gamma v$,
\begin{equation}
  D^2 - B^2 = \gamma^2 - \gamma^2v^2 = \gamma^2(1-v^2) = 1 \;\checkmark
\end{equation}
Ela não é falsa nem supérflua: é precisamente a equação que carrega a
informação que a simetria foi chamada para fornecer.

\paragraph{Ressalva sobre sinais.} O argumento de simetria, como enunciado em
\eqref{eq:reciproca}, fixa $D=+A$, mas \emph{não} fixa o sinal de $A$. A
solução $A=D=-\gamma$ também tem $\det=+1$ e também satisfaz tudo --- é o
boost composto com $PT$. Excluí-la exige um terceiro ingrediente:
preservação da orientação temporal, $\partial t'/\partial t>0$. Na prática,
o critério

\begin{center}\itshape
em $v\to0$ a transformação deve reduzir-se à identidade
\end{center}

\noindent
resolve os dois sinais de uma vez e é o mais econômico dos três argumentos:
ele seleciona a componente de $O(1,1)$ conexa à identidade, que é o grupo de
Lorentz próprio e ortocrono $SO^+(1,1)$ --- exatamente o conjunto dos boosts.

\vspace{10pt}
\noindent\textcolor{rgLinha}{\rule{\linewidth}{0.8pt}}
\vspace{4pt}

\noindent
{\small\color{rgGrafite}
\textbf{Resumo em uma linha.} A invariância do intervalo determina a
transformação a menos de reflexões discretas; a simetria $S\leftrightarrow
S'$ não fornece uma equação nova, apenas escolhe entre $D=+A$ e $D=-A$ ---
equivalentemente, entre $\det M=+1$ e $\det M=-1$.}

\end{document}
